% Part: sets-functions-relations
% Chapter: sets
% Section: unions-and-intersections

\documentclass[../../../include/open-logic-section]{subfiles}

\begin{document}

\olfileid[pt]{sfr}{set}{uni}
\olsection{Uniões e Interseções}

\begin{explain}
Na \olref[sfr][set][bas]{sec}, introduzimos definições de conjuntos por
abstração, isto é, definições da forma $\Setabs{x}{\phi(x)}$.
Aqui invocamos alguma propriedade~$\phi$, e essa propriedade pode mencionar
conjuntos que já definimos. Assim, por exemplo, se $A$ e~$B$ são conjuntos,
o conjunto $\Setabs{x}{x \in A \lor x \in B}$ consiste em todos aqueles
objetos que são !!{element}s de $A$ ou de~$B$, isto é, é o
conjunto que reúne os !!{element}s de $A$ e~$B$. Podemos visualizar
isto como na \olref{fig:union}, onde a área realçada indica os
!!{element}s dos dois conjuntos $A$ e~$B$ juntos.

\begin{figure}[H]
  \olasset{assets/diagrams/union.tikz}
  \caption{A união $A \cup B$ de dois conjuntos é o conjunto dos !!{element}s de
   $A$ juntamente com os de~$B$.}
  \ollabel{fig:union} 
\end{figure}

Esta operação sobre conjuntos---combiná-los---é muito útil e comum,
e por isso damos-lhe um nome formal e um símbolo. 
\end{explain}

\begin{defn}[União]
A \emph{união} de dois conjuntos $A$ e $B$, escrita $A \cup B$, é o
conjunto de todas as coisas que são !!{element}s de $A$, de $B$, ou de ambos.
\[
A \cup B = \Setabs{x}{x \in A \lor x \in B}
\]
\end{defn}

\begin{ex}
Como a multiplicidade de !!{element}s não importa, a união de dois
conjuntos que têm !!a{element} em comum contém esse !!{element} apenas uma vez,
por exemplo, $\{ a, b, c\} \cup \{ a, 0, 1\} = \{a, b, c, 0, 1\}$.

A união de um conjunto com um dos seus subconjuntos é simplesmente o conjunto maior: $\{a,
b, c \} \cup \{a \} = \{a, b, c\}$.

A união de um conjunto com o conjunto vazio é idêntica ao conjunto: $\{a,
b, c \} \cup \emptyset = \{a, b, c \}$.
\end{ex}

\begin{prob}
Prove que se $A \subseteq B$, então $A \cup B = B$.
\end{prob}

\begin{explain}
Também podemos considerar uma operação ``dual'' à união. É a
operação que forma o conjunto de todos os !!{element}s que são !!{element}s
de~$A$ e são também !!{element}s de~$B$. Esta operação chama-se
\emph{interseção}, e pode representar-se como na \olref{fig:intersection}.
\begin{figure}[H]
  \olasset{assets/diagrams/intersection.tikz}
  \caption{A interseção $A \cap B$ de dois conjuntos é o conjunto dos
    !!{element}s que têm em comum.}
  \ollabel{fig:intersection}
\end{figure}
\end{explain}

\begin{defn}[Interseção]
A \emph{interseção} de dois conjuntos $A$ e $B$, escrita $A \cap B$, é
o conjunto de todas as coisas que são !!{element}s de ambos $A$ e~$B$.
\[
A \cap B = \Setabs{x}{x \in A \land x \in B}
\]
Dois conjuntos dizem-se \emph{disjuntos} se a sua interseção é
vazia. Isso significa que não têm !!{element}s em comum.
\end{defn}

\begin{ex}
Se dois conjuntos não têm !!{element}s em comum, a sua interseção é vazia:
$\{ a, b, c\} \cap \{ 0, 1\} = \emptyset$.

Se dois conjuntos têm !!{element}s em comum, a sua interseção é o conjunto de
todos esses !!{element}s: $\{a, b, c \} \cap \{a, b, d \} = \{a, b\}$.

A interseção de um conjunto com um dos seus subconjuntos é simplesmente o conjunto menor:
$\{a, b, c\} \cap \{a, b\} = \{a, b\}$.

A interseção de qualquer conjunto com o conjunto vazio é vazia: $\{a, b, c \}
\cap \emptyset = \emptyset$.
\end{ex}

\begin{prob}
Prove rigorosamente que se $A \subseteq B$, então $A \cap B = A$.
\end{prob}

\begin{explain}
Também podemos formar a união ou a interseção de mais de dois
conjuntos. Uma forma elegante de tratar disso em geral é a
seguinte: suponhamos que reunimos num único conjunto todos os conjuntos dos quais queremos formar a união
(ou a interseção). Então podemos definir a união
de todos os nossos conjuntos originais como o conjunto de todos os objetos que pertencem a pelo
menos um !!{element} desse conjunto, e a interseção como o conjunto de
todos os objetos que pertencem a todo !!{element} desse conjunto.
\end{explain}

\begin{defn}
Se $A$ é um conjunto de conjuntos, então $\bigcup A$ é o conjunto dos !!{element}s dos
!!{element}s de~$A$:
\begin{align*}
\bigcup A & = \Setabs{x}{x \text{ pertence a !!a{element} de } A},
\text{ isto é,}\\
& = \Setabs{x}{\text{existe um } B \in A
  \text{ tal que } x \in B}
\end{align*}
\end{defn}

\begin{defn}
Se $A$ é um conjunto de conjuntos, então $\bigcap A$ é o conjunto dos objetos que
todos os elementos de~$A$ têm em comum:
\begin{align*}
\bigcap A & = \Setabs{x}{x \text{ pertence a todo !!{element} de } A},
\text{ isto é,}\\
 & = \Setabs{x}{\text{para todo } B \in A, x \in B}
\end{align*}
\end{defn}

\begin{ex}
Suponha $A = \{ \{ a, b \}, \{ a, d, e \}, \{ a, d \} \}$.
Então $\bigcup A = \{ a, b, d, e \}$ e $\bigcap A = \{ a \}$.
\end{ex}
\begin{prob}
	Mostre que se $A$ é um conjunto e $A \in B$, então $A \subseteq \bigcup B$.
\end{prob}

Também poderíamos fazer o mesmo para uma sequência de conjuntos $A_1$, $A_2$, \dots
\begin{align*}
\bigcup_i A_i & = \Setabs{x}{x \text{ pertence a um dos } A_i}\\
\bigcap_i A_i & = \Setabs{x}{x \text{ pertence a todos os } A_i}.
\end{align*}

Quando temos um \emph{índice} de conjuntos, isto é, algum conjunto $I$ tal que
consideramos $A_i$ para cada $i \in I$, podemos também usar estas
abreviaturas:
\begin{align*}
	\bigcup_{i \in I} A_i & = \bigcup \Setabs{A_i }{i \in I}\\
	\bigcap_{i \in I} A_i & = \bigcap\Setabs{A_i}{i \in I}
\end{align*}

Finalmente, podemos querer pensar no conjunto de todos os !!{element}s em~$A$
que não estão em~$B$. Podemos representá-lo como na \olref{difference}.

\begin{figure}[H]
  \olasset{assets/diagrams/difference.tikz}
  \caption{A diferença $A \setminus B$ de dois conjuntos é o conjunto dos
    !!{element}s de~$A$ que não são também !!{element}s de~$B$.}
  \ollabel{difference}
\end{figure}

\begin{defn}[Diferença]
A \emph{diferença} de conjuntos~$A \setminus B$ é o conjunto de todos os !!{element}s de
$A$ que não são também !!{element}s de~$B$, isto é,
\[
A\setminus B = \Setabs{x}{x\in A \text{ e } x \notin B}.
\]
\end{defn}

\begin{prob}
	Prove que se $A \subsetneq B$, então $B \setminus A \neq \emptyset$.
\end{prob}

\end{document}
